\section{Sparse stopping by encounter forests}\label{sec:sparse}
The four-row recursion will stop when a diagram is close to one row or column. We prove the needed bound directly from the paths of a small collection of cards. Branching expresses the trace for $s$ tracked cards as a binomial sum over first-row deficits. Inverting that sum cancels contributions supported on fewer than all the tracked cards. To preserve this cancellation, we couple all subsets of the tracked cards in one sweep.

Throughout this section, $n=2^d$, $T=T_n$, $B=TT^*$, and $B_\lambda=B_n(\lambda)$. Empty diagrams have dimension one.

\begin{lemma}[Sparse stopping]
\label{lem:sparse}
There is an absolute $N_0$ such that, if $n=2^d\ge N_0$ and
$1\le k\le n^{1/100}$, then
\[
 \sum_{\mu\vdash k}D_\mu\,
     \Tr B_{(n-k,\mu)}\le n^{-3k/10}.
\]
Consequently, for $n\ge N_0$, if
$k(\lambda)=n-\max\{\lambda_1,\lambda'_1\}\le n^{1/100}$,
then for every integer $r\ge4$,
\[
 D_\lambda\Tr B_\lambda^r\le1.
\]
There is a fixed even $r_0\ge4$ for which the latter assertion holds
for every power of two $n\ge2$ and every such $\lambda$.
\end{lemma}

\begin{proof}
\emph{Isolating the first-row deficit.}
Let $\mathcal I_s$ be the set of ordered injections from $[s]$ into
the $n$ positions, with $\mathcal I_0$ a singleton. The sweep induces
a transition matrix $T_s$ on this set. Write
\[
 a_s=\Tr_{\mathbb C^{\mathcal I_s}}(TT^*)
     =\sum_{x,y\in\mathcal I_s}T_s(x,y)^2.
\]
This is the ordinary, unnormalized trace on the indicated permutation module.

The permutation module on $\mathcal I_s$ is
$\operatorname{Ind}_{S_{n-s}}^{S_n}\mathbf1$. By repeated branching,
the multiplicity of $V_\lambda$ in it is the number of standard
tableaux of skew shape $\lambda/(n-s)$. Assume $s\le k$ and
$n\ge2k$. Every occurring partition has the form
$\lambda=(n-j,\mu)$ with $0\le j\le s$ and $\mu\vdash j$.
The skew shape consists of the tail $\mu$ and a row of $s-j$ cells
whose columns exceed $n-s\ge k\ge\mu_1$. These components share
neither a row nor a column. Choosing the $j$ entries of the tail and
then a standard tableau on it gives multiplicity
$\binom{s}{j}D_\mu$. Hence
\[
 a_s=\sum_{j=0}^s\binom{s}{j}b_j,
 \qquad
 b_j=\sum_{\mu\vdash j}D_\mu\,
                \Tr B_{(n-j,\mu)}.
\]
Binomial inversion gives
\begin{equation}
\label{eq:sparse-inverse}
 b_k=\sum_{s=0}^k(-1)^{k-s}\binom{k}{s}a_s\ge0,
\end{equation}
where positivity follows from the positive semidefiniteness of each
$B_\lambda$. It remains to bound this inverse without losing its
cancellation.

For fixed initial and final positions of one card, there is exactly
one possible path through the successive coordinate layers: at layer
$i$, its first $i$ processed coordinates have their final values and
its remaining coordinates still have their initial values. Thus fixed
endpoints for $s$ cards prescribe all their paths. If two tracked cards
use the same switch, the two prescribed coin requirements either
disagree, making the endpoints impossible, or agree, saving one coin.
If $e$ is the number of switches shared by two tracked cards along a
realized sweep, the transition probability to its realized endpoint is
therefore $2^{-ds+e}=n^{-s}2^e$. Averaging first over the realized
endpoint and then over a uniform initial injection gives the exact
identity
\begin{equation}
\label{eq:sparse-trace}
 a_s=\frac{(n)_s}{n^s}\,\mathbb E\,2^e,
 \qquad (n)_s=n(n-1)\cdots(n-s+1).
\end{equation}

Choose $k$ distinct uniform initial positions $X_1,\ldots,X_k$,
independently of the coins of one complete sweep. For this realized
sweep, form a graph $G$ on all $n$ initial positions, joining two
positions if their cards share a switch. Each vertex has degree $d$.
The graph is simple: after two cards meet at coordinate $i$, their
positions differ at that coordinate forever, so they cannot meet at a
later coordinate. Let $G_X$ be the graph induced on the selected tags,
and let $e(S)$ count its edges with both endpoints in $S\subseteq[k]$.
The marginal initial positions of every such subset are a uniform
injection. Thus, with
\[
 w(S)=\frac{(n)_{|S|}}{n^{|S|}}
     =\prod_{i\in S}\left(1-
             \frac{\#\{j\in S:j<i\}}n\right),
\]
Equations~\eqref{eq:sparse-trace} and \eqref{eq:sparse-inverse} imply
\begin{equation}
\label{eq:sparse-coupled}
 b_k=\mathbb E\sum_{S\subseteq[k]}
             (-1)^{k-|S|}w(S)2^{e(S)}.
\end{equation}

A predecessor selection is a collection $D$ of directed pairs
$(i,j)$ with $1\le j<i\le k$, at most one pair having any given
first coordinate. Write $|D|$ for its number of pairs and
$\operatorname{supp}D$ for their endpoints. Expanding the product
defining $w(S)$ gives
\[
 w(S)=\sum_{D:\operatorname{supp}D\subseteq S}
                         (-n^{-1})^{|D|}.
\]
Also $2^{e(S)}=\sum_{H\subseteq E(G_X[S])}1$. If
$\operatorname{supp}H$ denotes the vertices incident to the edges
of $H$, then summing over $S$ in \eqref{eq:sparse-coupled} yields the
exact cancellation identity
\begin{equation}
\label{eq:sparse-cover}
 b_k=\mathbb E\sum_{H\subseteq E(G_X)}\ 
       \sum_{D:\operatorname{supp}H\cup\operatorname{supp}D=[k]}
                     (-n^{-1})^{|D|}.
\end{equation}
Indeed, a term supported on $U$ has coefficient
$\sum_{S\supseteq U}(-1)^{k-|S|}$, which is zero unless $U=[k]$.
Thus only an encounter-edge set and a predecessor selection that
together cover every tag survive the inversion. We may now bound their
absolute sum without losing that covering constraint.

\medskip\noindent
\emph{Counting the covering configurations.}
Two elementary facts about $G$ control the possible encounter
subgraphs. First, every set of $v$ vertices spans at most
\begin{equation}
\label{eq:sparse-edge-count}
 \frac v2\log_2v
\end{equation}
edges, with $0\log_2 0=0$. To prove this by induction on $d$, split
the initial positions according to the last processed coordinate.
Before the last layer, these two halves remain separate. If the
chosen set has $v_0,v_1$ positions in them, induction bounds the
earlier encounters by
$\tfrac12v_0\log_2v_0+\tfrac12v_1\log_2v_1$, and the last matching
adds at most $\min(v_0,v_1)$ edges. This is at most
$\tfrac12v\log_2v$ because the binary entropy
$H_2(p)=-p\log_2p-(1-p)\log_2(1-p)$, with $0\log_2 0=0$, satisfies
$H_2(p)\ge2\min(p,1-p)$ by concavity on each half of $[0,1]$.

Second, if $F$ is a specified forest on $v$ selected tags with $c$
components, then, conditional on the entire sweep,
\begin{equation}
\label{eq:sparse-forest-probability}
 \Pr(F\subseteq G_X\mid G)
 \le\frac{n^cd^{v-c}}{(n)_v}
 \le2\left(\frac dn\right)^{v-c}.
\end{equation}
Choose one root position per component and then successively place
each child at one of at most $d$ neighbors of its parent. This gives
the first inequality even if some counted placements fail to be
injective. For $v\le k\le n^{1/100}$ and sufficiently large $n$,
\[
 \frac{(n)_v}{n^v}
 =\prod_{i=0}^{v-1}(1-i/n)
 \ge1-\frac{v(v-1)}{2n}\ge\frac12,
\]
which gives the second inequality.

Let an encounter-edge set $H$ use $v$ vertices and have $c$ nontrivial
connected components. Then $2c\le v$, and $H$ contains a spanning
forest with $v-c$ edges. For fixed $v,c$, there are at most
$(2k)^{2v}$ choices for its vertex set and forest: choose the vertex
set, root every component at its least vertex, and record a parent
or a root marker for every vertex. For any realization and any such
forest, the number of possible edge sets extending it is at most
\[
 2^{e(G_X[V])}\le v^{v/2}\le k^{v/2},
\]
by \eqref{eq:sparse-edge-count}. A canonical choice of spanning
forest for each $H$, or simply overcounting all spanning forests,
therefore bounds the expected number of these $H$ by
\[
 2(2k)^{2v}k^{v/2}(d/n)^{v-c}.
\]
For $v=c=0$, the empty edge set contributes one, which also obeys
this upper bound.

There are at most $\binom{k}{q}k^q\le(2k)^{2q}$ predecessor
selections with $q$ pairs. The covering condition in
\eqref{eq:sparse-cover} implies $v+2q\ge k$. Since $d/n\le1$,
$v-c\ge v/2$, and there are at most $(k+1)^3$ triples $(v,c,q)$,
the preceding estimates show that
\begin{equation}
\label{eq:sparse-sum}
 b_k\le 2\sum_{\substack{0\le v,c,q\le k\\
                         2c\le v,\ v+2q\ge k}}
       \left(4k^{5/2}\sqrt{d/n}\right)^v
       \left(4k^2/n\right)^q.
\end{equation}
Impossible triples only enlarge this upper bound. Uniformly for
$1\le k\le n^{1/100}$, sufficiently large $n$ satisfies
\[
 4k^{5/2}\sqrt{d/n}\le n^{-2/5},\qquad
 2k/\sqrt n\le n^{-2/5},\qquad
 2(k+1)^3\le n^{k/10}.
\]
For example, the first left side is at most
$4\sqrt{\log_2n}\,n^{-19/40}$, and the last left side is at most
$16n^{3/100}$. Every term in \eqref{eq:sparse-sum} is therefore at
most $n^{-2(v+2q)/5}\le n^{-2k/5}$, proving
$0\le b_k\le n^{-3k/10}$.

\medskip\noindent
\emph{From the trace bound to stopping.}
If $\lambda_1=n-k$ with $1\le k\le n^{1/100}$, its term in $b_k$
has coefficient $D_\mu\ge1$, so
$\Tr B_\lambda\le n^{-3k/10}$. The eigenvalues are
nonnegative, and consequently
\[
 \Tr B_\lambda^r
 \le(\Tr B_\lambda)^r.
\]
Choosing the labels of the $k$ cells below the first row and then
their order gives $D_\lambda\le\binom nk k!\le n^k$. Hence
\[
 D_\lambda\Tr B_\lambda^r
 \le n^{(1-3r/10)k}\le1\qquad(r\ge4).
\]

If instead $\lambda'_1=n-k>n/2$, every matching projection vanishes
on $V_\lambda$. Indeed, its fixed-space dimension is the
multiplicity of $V_\lambda$ in
$\operatorname{Ind}_{S_2^{n/2}}^{S_n}\mathbf1$. Repeated horizontal
Pieri shows that every constituent is obtained by adding $n/2$
horizontal strips of size two. Each strip adds at most one cell to
the first column, so every constituent has first-column height at
most $n/2$. Thus $B_\lambda=0$ in the case under consideration.
For $k=0$, the trivial representation has the claimed value one
and the sign representation is annihilated by a matching projection.

Finally, Lemma~\ref{lem:gap} gives operator norm strictly less than
one on every nontrivial irreducible at each fixed $n$.
There are only finitely many powers of two below $N_0$ and finitely
many irreducibles at each of them. Increasing $r$ to one fixed even
$r_0\ge4$ therefore gives the stated stopping estimate also at these
remaining sizes. Larger powers preserve the estimate because $B$
is a positive semidefinite contraction.
\end{proof}
